%%%%%%%%%%%%%%%%%%%%%%%%%%%%%%%%%%%%%%%%%%%%%%%%%%%%%%%%%%%%%%%%%%%%%
%% Artifact Appendix for mlir-opt-omp (AE template V20190108)
%%
%% Paste the part between \appendix and \subsection{Methodology}
%% into the paper.  Items marked "TODO" are the ones only you can fill
%% (DOI, license, exact machine, measured wall-clock times).
%%%%%%%%%%%%%%%%%%%%%%%%%%%%%%%%%%%%%%%%%%%%%%%%%%%%%%%%%%%%%%%%%%%%%

\appendix
\section{Artifact Appendix}

%%%%%%%%%%%%%%%%%%%%%%%%%%%%%%%%%%%%%%%%%%%%%%%%%%%%%%%%%%%%%%%%%%%%%
\subsection{Abstract}

The artifact is \texttt{mlir-opt-omp}, an out-of-tree \texttt{mlir-opt} that
lowers the MLIR OpenMP dialect (\texttt{omp.*}) to calls into an OpenMP
runtime library, driven by a declarative rule file (\texttt{rules.dsl}) that
describes each runtime's ABI.  It contains the compiler sources, the rule file
for the three runtimes evaluated in the paper (Intel/LLVM \texttt{\_\_kmpc\_*},
GCC \texttt{GOMP\_*} and the PMSIS cluster API of PULP/GAP8), a lit+FileCheck
regression suite over the passes, and shell drivers that take PolyBench/C
kernels end to end (C $\rightarrow$ ClangIR $\rightarrow$ \texttt{mlir-opt-omp}
$\rightarrow$ binary) and compare them against the same kernels built by a
stock OpenMP compiler, for both correctness and performance.

Reviewers can reproduce the IR-level claims --- that one rule file, and only
the rule file, decides which runtime a construct is lowered onto --- in minutes
on any machine where LLVM/MLIR builds.  Reproducing the end-to-end correctness
and speedup numbers additionally needs the ClangIR front-end and a host OpenMP
runtime; reproducing the PULP/GAP8 numbers needs the GAP SDK with the
\texttt{gvsoc} simulator (no board is required).

\subsection{Artifact check-list (meta-information)}

{\small
\begin{itemize}
  \item {\bf Algorithm: } rule-driven lowering of OpenMP constructs: a
    declarative per-runtime ABI description is evaluated against each
    \texttt{omp.*} operation into a lowering plan, realised by three MLIR
    passes (rule evaluation, outlining + work-sharing loop lowering, call
    emission); plus an optional redundant team-barrier elimination pass.
  \item {\bf Program: } PolyBench/C 4.2.1 OpenMP kernels (\texttt{gemm} and
    \texttt{atax} vendored in the artifact; the full suite from an external
    checkout), and 74 hand-written MLIR regression tests.
  \item {\bf Compilation: } CMake $\geq$ 3.20, C++17, Ninja; LLVM/MLIR 23 built
    with ClangIR enabled (EEESlab fork, branch
    \texttt{users/lucap/cir-omp-clauses}, commit \texttt{9be70f177d9b}).
    Baseline kernels are compiled with stock \texttt{clang} (\texttt{iomp}),
    \texttt{gcc} (\texttt{libgomp}) or the PULP-SDK \texttt{gcc}
    (\texttt{pmsis}), all at \texttt{-O3} with matched strict-FP flags.
  \item {\bf Transformations: } \texttt{--omp-to-omp-lower},
    \texttt{--omp-outline}, \texttt{--omp-lower-plan}, and optionally
    \texttt{--omp-barrier-elim}, selected per runtime with
    \texttt{--omp-lower-runtime=\{iomp,libgomp,pmsis\}}.
  \item {\bf Binary: } the \texttt{mlir-opt-omp} tool, plus the per-kernel host
    binaries (and, for \texttt{pmsis}, RISC-V ELFs) that the drivers build from
    source at run time.
  \item {\bf Data set: } PolyBench's own generated inputs, selected by dataset
    macro: \texttt{MINI\_DATASET} for correctness,
    \texttt{LARGE\_DATASET} for host performance, \texttt{MINI\_DATASET} for
    PULP/GAP8 (memory bound).
  \item {\bf Run-time environment: } Linux x86\_64 (Ubuntu 22.04/24.04) with
    \texttt{libomp} or \texttt{libgomp}; runs use
    \texttt{OMP\_PROC\_BIND=true}, \texttt{OMP\_PLACES=cores} and a spinning
    wait policy (\texttt{OMP\_WAIT\_POLICY=ACTIVE},
    \texttt{KMP\_BLOCKTIME=infinite}, \texttt{GOMP\_SPINCOUNT=infinite}) so that
    timings are repeatable.  The \texttt{pmsis} target additionally needs the
    GAP SDK and the \texttt{gvsoc} simulator.
  \item {\bf Hardware: } a multi-core x86\_64 machine.  The reported host
    numbers were taken on TODO (CPU model, core count), using every hardware
    thread (\texttt{THREADS=16}) on an otherwise idle machine.  The
    PULP/GAP8 results are obtained on the cycle-accurate \texttt{gvsoc}
    simulator; no hardware board is needed.
  \item {\bf Run-time state: } threads pinned to cores and kept spinning; the
    machine must be otherwise idle, since idle workers hold their core.
  \item {\bf Execution: } POSIX shell drivers; each timing cell is run
    \texttt{REPS}=10 times, minimum and maximum dropped, mean $\pm$ standard
    deviation reported, cells with relative std-dev above 5\% flagged as noisy.
    On \texttt{gvsoc} each cell is run once (the simulator is deterministic).
  \item {\bf Metrics: } correctness --- bit-identical PolyBench array dumps
    between the reference build and ours.  Performance --- cycles from
    PolyBench's TSC timer, reported as the self-relative parallel speedup of
    each toolchain (\texttt{speedup\_native}, \texttt{speedup\_opt}) and as our
    parallel/sequential time against the native one
    (\texttt{opt\_vs\_native\_par}, \texttt{opt\_vs\_native\_seq}), summarised
    over the suite with the geometric mean.  On PULP additionally the size of
    the linked ELF.
  \item {\bf Output: } console log plus, under
    \texttt{test/Integration/results/<runtime>/},
    \texttt{results\_correctness.csv} (\texttt{kernel;PASS|FAIL|ERROR}),
    \texttt{results\_performance.csv} (per-kernel rows and a \texttt{GEOMEAN}
    summary row) and, with \texttt{PLOT=true}, the speedup bar chart used for
    Figures TODO (and the binary-size chart on \texttt{pmsis}); the intermediate
    LLVM IR of every build is kept for inspection.
  \item {\bf Experiments: } (i) the regression suite (\texttt{check-omp});
    (ii) \texttt{run\_correctness.sh}; (iii) \texttt{run\_performance.sh};
    (iv) \texttt{tasks/run\_tasks.sh} for the task construct; each of (ii)--(iv)
    per runtime.  Driven by environment variables, no manual editing.
    Expected variation: the correctness results are exact and deterministic;
    speedups vary with the machine, run-to-run noise is bounded by the
    repetition scheme above.
  \item {\bf How much disk space required (approximately)?: } about 1\,GB for
    the artifact and its results; about 80\,GB more for a Release LLVM/MLIR +
    ClangIR build tree, if one is not already available.
  \item {\bf How much time is needed to prepare workflow (approximately)?: }
    about 15 minutes given an existing LLVM/MLIR+ClangIR build (configure,
    build the tool, fill in \texttt{local.env}); 2--4 hours otherwise,
    essentially all of it the LLVM build.
  \item {\bf How much time is needed to complete experiments (approximately)?: }
    under a minute for the regression suite, a few minutes for the bundled
    kernels, and about TODO hours for the full PolyBench suite at
    \texttt{LARGE\_DATASET} on both host runtimes; the \texttt{gvsoc} runs take
    about TODO hours.
  \item {\bf Publicly available?: } yes ---
    \url{https://github.com/EEESlab/mlir-opt-omp}, together with the ClangIR
    fork at \url{https://github.com/EEESlab/llvm-project}.
  \item {\bf Code licenses (if publicly available)?: } TODO (the repository
    carries no \texttt{LICENSE} file yet; Apache-2.0 with LLVM exceptions is
    the natural choice for an out-of-tree MLIR tool).
  \item {\bf Data licenses (if publicly available)?: } PolyBench/C 4.2.1 and
    its OpenMP port, under their own (BSD-style) license.
  \item {\bf Workflow framework used?: } none --- CMake, \texttt{lit} and
    POSIX shell drivers configured through two environment files.
  \item {\bf Archived (provide DOI)?: } TODO (Zenodo DOI of the tagged
    snapshot).
\end{itemize}
}

%%%%%%%%%%%%%%%%%%%%%%%%%%%%%%%%%%%%%%%%%%%%%%%%%%%%%%%%%%%%%%%%%%%%%
\subsection{Description}

\subsubsection{How delivered}

As a public Git repository, \url{https://github.com/EEESlab/mlir-opt-omp},
and as an archived snapshot of the tag used for the paper (TODO: DOI).  The
snapshot holds the compiler sources, \texttt{rules.dsl}, the regression suite,
the integration drivers and the two vendored PolyBench kernels --- roughly
TODO\,MB.  The ClangIR front-end is a separate repository (a fork of
\texttt{llvm-project}) and is not vendored; the exact commit it is tested
against is pinned in the top-level \texttt{README.md}.

\subsubsection{Hardware dependencies}

A multi-core x86\_64 Linux machine for the \texttt{iomp} and \texttt{libgomp}
experiments; the parallel runs pin one thread per physical core, so the core
count sets the achievable speedup.  Building LLVM/MLIR with ClangIR is the
demanding step (a machine with $\geq$16\,GB of RAM and several cores is
strongly recommended).  The PULP/GAP8 experiments need no hardware: they run on
the \texttt{gvsoc} simulator.

\subsubsection{Software dependencies}

CMake $\geq$ 3.20, a C++17 compiler and Ninja; an LLVM/MLIR 23 build ---
with ClangIR enabled if the C-to-MLIR path is to be exercised, any MLIR 23
build or install otherwise --- configured with \texttt{LLVM\_INSTALL\_UTILS=ON}
so that \texttt{FileCheck} and \texttt{lit} are available for the regression
suite.  The end-to-end drivers additionally need a stock OpenMP compiler
(\texttt{clang} for \texttt{iomp}, \texttt{gcc} for \texttt{libgomp}) and its
OpenMP headers, and \texttt{python3} with \texttt{matplotlib} and
\texttt{numpy} if the figures are to be rendered (the runs succeed without
them, only the plot is skipped).  For \texttt{pmsis}: the GAP SDK with
\texttt{gvsoc}, the GAP RISC-V GCC toolchain, an LLVM build with the RISC-V
(\texttt{xpulpv}) back-end, and the external PolyBench-PULP harness that
\texttt{PULP\_APP\_DIR} points at.

\subsubsection{Data sets}

PolyBench/C 4.2.1 with OpenMP pragmas.  Two kernels (\texttt{gemm} and
\texttt{atax}, plus the PolyBench support code) are vendored under
\texttt{test/Integration/kernels/} so that a first run is self-contained; the
full suite is used by pointing \texttt{POLYBENCH} at an external checkout and
setting \texttt{SUITE=full}.  Inputs are generated by PolyBench itself from the
selected dataset macro; no external data files are needed.

%%%%%%%%%%%%%%%%%%%%%%%%%%%%%%%%%%%%%%%%%%%%%%%%%%%%%%%%%%%%%%%%%%%%%
\subsection{Installation}

\noindent 1.\ Obtain an LLVM/MLIR 23 build with ClangIR:

{\small
\begin{verbatim}
git clone -b users/lucap/cir-omp-clauses \
  https://github.com/EEESlab/llvm-project.git
git -C llvm-project checkout 9be70f177d9b
cmake -S llvm-project/llvm -B llvm-project/build -G Ninja \
  -DCMAKE_BUILD_TYPE=Release \
  -DLLVM_ENABLE_PROJECTS="clang;mlir;lld" \
  -DCLANG_ENABLE_CIR=ON \
  -DLLVM_TARGETS_TO_BUILD="X86;RISCV" \
  -DLLVM_INSTALL_UTILS=ON
cmake --build llvm-project/build
\end{verbatim}
}

\noindent 2.\ Build the tool against it:

{\small
\begin{verbatim}
git clone https://github.com/EEESlab/mlir-opt-omp.git
cmake -S mlir-opt-omp -B mlir-opt-omp/build -G Ninja \
  -DMLIR_DIR=<LLVM_DIR>/lib/cmake/mlir
cmake --build mlir-opt-omp/build
\end{verbatim}
}

\noindent where \texttt{<LLVM\_DIR>} is either the build tree from step 1 or an
install prefix.

\noindent 3.\ Tell the end-to-end drivers where the tools are (a single file at
the repository root; the drivers validate it at start-up and report which entry
is wrong):

{\small
\begin{verbatim}
cp local.env.example local.env
$EDITOR local.env     # LLVM_BIN, OMP_TOOL_BIN, INC_OMP
\end{verbatim}
}

\noindent Ready-made configurations for the machines the work was developed on
are under \texttt{docs/setups/}.

%%%%%%%%%%%%%%%%%%%%%%%%%%%%%%%%%%%%%%%%%%%%%%%%%%%%%%%%%%%%%%%%%%%%%
\subsection{Experiment workflow}

\noindent {\bf E1 --- passes (minutes, no runtime needed).}

{\small
\begin{verbatim}
cmake --build build --target check-omp
\end{verbatim}
}

\noindent The same module lowered onto two runtimes, which is the claim the
rule file embodies:

{\small
\begin{verbatim}
build/mlir-opt-omp \
  test/Regression/wsloop/schedule-static/\
schedule-static-iomp.mlir \
  --omp-lower-dsl=rules.dsl --omp-to-omp-lower \
  --omp-outline --omp-lower-plan \
  --omp-lower-runtime=iomp | grep 'call @'
# then again with --omp-lower-runtime=libgomp
\end{verbatim}
}

\noindent {\bf E2 --- end-to-end correctness.}

{\small
\begin{verbatim}
cd test/Integration
./run_correctness.sh                    # iomp, bundled
RUNTIME=libgomp ./run_correctness.sh
SUITE=full POLYBENCH=/path/to/checkout \
  ./run_correctness.sh
\end{verbatim}
}

\noindent {\bf E3 --- performance (Figures TODO).}

{\small
\begin{verbatim}
POLYBENCH=/path/to/checkout \
  RUN_ENV=configs/paper-libgomp.env PLOT=true ./run_performance.sh
POLYBENCH=/path/to/checkout \
  RUN_ENV=configs/paper-iomp.env    PLOT=true ./run_performance.sh

# Shorter, less confident: an environment variable wins over the config file.
#   REPS=5  halves the time, keeps a deviation over three samples
#   REPS=3  3.3x faster, the median of three, no deviation
# The dataset is deliberately not the knob to turn: fewer repetitions cost
# confidence in a value, a smaller dataset changes which value is measured.
\end{verbatim}
}

\noindent {\bf E4 --- tasks.}
\texttt{tasks/run\_tasks.sh libgomp} (and \texttt{iomp}), two end-to-end checks
against the real runtime library.

\noindent {\bf E5 --- PULP/GAP8 (Figure TODO), only with the GAP SDK
installed.}

{\small
\begin{verbatim}
RUNTIME=pmsis ./run_correctness.sh
RUN_ENV=configs/paper-pmsis.env PLOT=true ./run_performance.sh

# PLOT=true renders both PULP figures here: the speedup chart, and a second
# ..._size.png from the size_* columns that only this target writes.
# There is no shorter variant: gvsoc is deterministic so every cell runs once,
# and DATASET is already pinned to MINI by the GAP8 memory budget.
\end{verbatim}
}

%%%%%%%%%%%%%%%%%%%%%%%%%%%%%%%%%%%%%%%%%%%%%%%%%%%%%%%%%%%%%%%%%%%%%
\subsection{Evaluation and expected result}

E1 must report all 74 regression tests passing, and the two invocations must
show the same loop lowered into \texttt{\_\_kmpc\_*} calls and into
\texttt{GOMP\_*} calls respectively --- the only difference between the two
commands being the runtime name.

E2 must report \texttt{PASS} for every kernel in
\texttt{results\_correctness.csv}: the array dumps of our binary and of the
stock-compiler binary are compared bit for bit, so this is exact and
deterministic, and the driver exits non-zero if any kernel is not \texttt{PASS}.

E3 writes \texttt{results\_performance.csv} with one row per kernel and a
\texttt{GEOMEAN} row.  The claim to check is the \texttt{opt\_vs\_native\_par}
column and its geometric mean: our lowering is expected to stay within TODO of
the native toolchain, and \texttt{speedup\_opt} is expected to track
\texttt{speedup\_native} kernel by kernel (this is what the bar chart shows,
and it corresponds to Figures TODO of the paper).  Absolute times, and hence
individual speedups, depend on the machine; the shape of the comparison should
not.  Cells flagged as noisy (relative std-dev above
\texttt{VARIANCE\_ACCEPTED}) indicate the machine was not idle.

E5 reproduces the PULP/GAP8 columns: cycles scraped from the \texttt{gvsoc}
console and the size of the linked ELF.  Being a deterministic simulation,
these numbers should match the paper's exactly for the same SDK version.

%%%%%%%%%%%%%%%%%%%%%%%%%%%%%%%%%%%%%%%%%%%%%%%%%%%%%%%%%%%%%%%%%%%%%
\subsection{Experiment customization}

Everything is driven by environment variables, set inline, in
\texttt{test/Integration/run.env} or in \texttt{local.env}, in that order of
precedence: \texttt{RUNTIME} (\texttt{iomp}, \texttt{libgomp},
\texttt{pmsis}), \texttt{SUITE} (\texttt{bundled}, \texttt{full}) or an
explicit \texttt{KERNELS} list, \texttt{DATASET}, \texttt{THREADS},
\texttt{REPS}, \texttt{VARIANCE\_ACCEPTED}, \texttt{BARRIER\_ELIM},
\texttt{PLOT} and \texttt{OUTDIR}.  Results of different runtimes and of the
barrier-elimination variant land in separate directories, so runs never
overwrite each other.  The full reference is the ``Configuration reference''
section of \texttt{test/Integration/README.md}.

The lowering itself is customised without rebuilding: \texttt{rules.dsl} is
read at run time, so adding a runtime or changing how a construct is emitted is
an edit to that file, passed with \texttt{--omp-lower-dsl=<file>} (or
\texttt{RULES=<file>} for the drivers).  Constructs whose lowering needs
primitives the DSL does not express still require C++ changes in the outlining
pass.

%%%%%%%%%%%%%%%%%%%%%%%%%%%%%%%%%%%%%%%%%%%%%%%%%%%%%%%%%%%%%%%%%%%%%
\subsection{Notes}

The OpenMP emission this pipeline consumes is not upstream in ClangIR yet,
which is why the pinned fork is required for the C-to-MLIR path; starting from
\texttt{.mlir} input, any LLVM/MLIR 23 build works, and E1 is reproducible
there.  \texttt{INC\_OMP} must point at the OpenMP headers of the same GCC the
\texttt{libgomp} baseline uses, or the front-end and the baseline disagree.
The strict floating-point flags applied to both sides are deliberate: without
them FMA contraction and reordered reductions make iterative kernels diverge
even though both binaries are IEEE-correct.  The spinning wait policy is there
for repeatability, not speed --- numbers taken with it are not comparable with
numbers taken without it.  Finally, the \texttt{pmsis} path drives an external
PolyBench-PULP harness (\texttt{PULP\_APP\_DIR}) that is not part of this
repository; TODO: say where reviewers get it.
